\section{Síntesis y ampliación}

Esta sección condensa el episodio completo en seis conclusiones. Primera: lo decisivo en la revolución robótica es el paso de los modelos especializados a los modelos generales. Segunda: la vía inicial consistió en usar LLM/VLM/Code LM para transformar Planning, Perception y Actuation. Tercera: un VLA propiamente dicho integra Vision, Language y Action en un modelo de extremo a extremo. Cuarta: las líneas como RT-2/OpenVLA transfieren el conocimiento de internet mediante un VLM, pero el control de action a alta frecuencia todavía requiere un tratamiento específico. Quinta: Diffusion Policy/RDT muestran que la propia acción puede modelarse con modelos generativos. Sexta: las direcciones futuras son los modelos del mundo, la unificación de comprensión y predicción, y el aprendizaje por refuerzo en línea.

\begin{importantbox}{Ubicar la clase de VLA con pantalla compartida dentro de la serie de IA/internet de Zhang Xiaojun (张小珺)}
Este episodio es la base técnica de las entrevistas sobre inteligencia corporizada como EP106, EP109 y EP121. No es una historia de emprendimiento, sino una hoja de ruta de los modelos base para robots, y explica por qué VLA se volvió un término central en la discusión sobre robótica de 2025.
\end{importantbox}

\subsection{Tabla de consulta rápida de artículos clásicos}

\noindent\begin{tabularx}{\textwidth}{p{0.20\textwidth}p{0.30\textwidth}X}
\toprule
Artículo/sistema & Aportación principal & Lugar en esta clase \\
\midrule
SayCan & Planificación con LLM + viabilidad por affordance & Ejemplo representativo de sustituir Planning con un modelo de lenguaje. \\
Inner Monologue & Retroalimentación del entorno y autocorrección expresada en lenguaje & Hace que la retroalimentación de la ejecución vuelva a la planificación. \\
DoReMi & Detección y recuperación de desajustes entre el plan y la ejecución & Atiende las fallas de ejecución en robots reales. \\
VoxPoser & 3D value maps & Convierte objetivos en lenguaje en campos de operación espacial. \\
ALOHA & Manipulación bimanual de bajo costo y ACT & Datos de demostración y aprendizaje de secuencias de acciones. \\
RT-1/RT-2 & Transformer para robots / transferencia VLA & Línea principal de Google Robotics. \\
RT-X/OpenVLA & Datos de múltiples cuerpos robóticos y VLA de código abierto & Datos e infraestructura de código abierto. \\
HiRT/pi0/GR00T N1 & VLM + Action Head/Expert & Tratan la action con más rigor. \\
Diffusion Policy/RDT & Difusión de acciones y modelo base de manipulación bimanual & action policy generativa. \\
UP-VLA/iRe-VLA & Unificación de comprensión y predicción, y RL en línea & Direcciones futuras. \\
\bottomrule
\end{tabularx}

\subsection{Preguntas para el seguimiento}

Por último, volvemos al marco de observación a largo plazo de esta serie de entrevistas de IA/internet de Zhang Xiaojun: VLA no es el nombre de un modelo aislado, sino un conjunto de problemas de sistema en torno a los datos, los modelos, el hardware, la adopción comercial y el entrenamiento seguro. Las siguientes preguntas sirven como lista de seguimiento para juzgar si los modelos base para robots han entrado de verdad en una etapa de scale comparable a la de los grandes modelos de lenguaje.

\begin{enumerate}
\item ¿Aparecerá la scaling law de los VLA de manera tan estable como en los modelos de lenguaje?
\item ¿Podrán los conjuntos de datos de múltiples cuerpos robóticos ser lo bastante grandes y estandarizados para sostener políticas verdaderamente generales?
\item Entre un VLM backbone que emite la action directamente y un VLM + Action Head, ¿cuál de las dos líneas es más escalable?
\item ¿Se convertirán los modelos del mundo en el factor decisivo para la generalización robótica, o bastará con la propia action policy?
\item ¿Puede el aprendizaje por refuerzo en línea ejecutarse de forma segura, eficiente y sostenible en robots reales?
\item ¿Los robots humanoides llegarán primero a la aplicación práctica gracias al cuerpo robótico y la teleoperación, o tendrán que esperar a que maduren las capacidades de los VLA?
\end{enumerate}

\subsection{Lecturas adicionales}

\begin{itemize}
\item Si le interesa la historia académica de la inteligencia corporizada, puede compararse con la entrevista a Wang He (王鹤) en EP106.
\item Si le interesan la simulación y los datos sintéticos, puede compararse con la entrevista a Xie Chen (谢晨) de Guanglun (光轮) en EP109.
\item Si le interesan los modelos del mundo multimodales y el razonamiento visual, puede compararse con la entrevista a Zhang Xiangyu (张祥雨) en EP102.
\item Si le interesan la industrialización de la robótica y las vías de entrada por hardware, puede compararse con las entrevistas a Rokid en EP104 y a Tan Jie (谭捷) en EP121.
\end{itemize}

\end{document}
